\documentclass[10pt,a4paper]{amsart}
\usepackage[margin=19mm]{geometry}
\usepackage{amsmath,amssymb,mathtools}
\usepackage[scr=rsfs,cal=euler]{mathalfa}
\usepackage{xfrac}
\usepackage[unicode,hidelinks,bookmarks=false]{hyperref}
\newcommand{\ie}{i.e.\ }
\newcommand{\cf}{cf.\ }
\newcommand{\resp}{respectively\ }
\newcommand{\Subsection}{\subsection}
\providecommand{\nobreakdash}{\nobreak\hbox{-}}
\setlength{\emergencystretch}{2em}
\multlinegap0pt
\usepackage{amssymb}
\usepackage{dsfont}
\usepackage{mathtools}
\usepackage{xcolor}
\usepackage[shortlabels]{enumitem}
\newtheorem{lemma}{Lemma}
\newtheorem{theorem}{Theorem}
\newcommand{\M}{\mathcal M}
\newcommand{\ideal}[1]{\langle #1\rangle}
\title{Maximal-Hull $z$-Ideals, Congruence Closures, and Coherent Frames of Commutative Semirings}
\author{Pubali Sengupta, Amartya Goswami, Pronay Biswas, Sujit Kumar Sardar$^\dagger$}
\date{Source: arXiv:2607.07319v1 (CC BY 4.0)}
\begin{document}
\maketitle

\noindent\emph{Source and licence.} Maximal-Hull $z$-Ideals, Congruence Closures, and Coherent Frames of Commutative Semirings. Version arXiv:2607.07319v1 (CC BY 4.0), distributed by arXiv under the Creative Commons Attribution 4.0 International licence, http://creativecommons.org/licenses/by/4.0/, as recorded in the arXiv licence metadata for this exact version and shown on the abstract page https://arxiv.org/abs/2607.07319v1. Full licence notice: \texttt{licenses/arxiv-2607.07319-CC-BY-4.0.txt}.\par\smallskip

\noindent\textit{Editorial scope.} Counted unit: the article's theorem thm:regularity (source0.tex lines 785--794). The article's complete original proof of it is delivered with it (source0.tex lines 795--827, one \texttt{proof} environment of 33 lines). No mathematics is written, repaired or re-derived: the statement and the proof are the article's own lines, in the article's own notation, and no retained line is substituted or re-flowed.  Retained uncounted as the source-local context the proof needs: lines 547--553.  Nothing was omitted from any retained span, so the package carries no omission record.  No retained line contains a citation, so the wrapper carries no bibliography and no bibliography entry.  No retained line contains a figure environment, an image-inclusion command or drawing code, so no image is delivered and the package declares no asset dependency.  Editorial formatting only, in addition: an AMS \texttt{amsart} preamble that restates the article's own declarations and macros used by the retained lines, namely the environments \texttt{lemma}, \texttt{theorem} on separate counters, together with the standard packages those lines need.  The retained environments print in this document as Lemma 1, Theorem 1 in order of appearance, whereas the article prints the same items with the numbering of its own full document.  The complete native source of the article as distributed in the arXiv e-print for this exact version is bundled unchanged as \texttt{sources/204696/source0.tex}.
\begin{lemma}\label{loc6}
For all $a,b\in S$,
\begin{enumerate}
\item $\M(ab)=\M(a)\cup\M(b)$;
\item $\M(a)\cap\M(b)\subseteq\M(a+b)$.
\end{enumerate}
\end{lemma}

\begin{theorem}\label{thm:regularity}
Let $S$ be a semiring such that $E(S)=\mathrm{Comp}(S)$.
Then the following are equivalent.
\begin{enumerate}
\item $S$ is von Neumann regular.
\item Every ideal of $S$ is a strong $z$-ideal.
\item Every ideal of $S$ is a $z$-ideal.
\item Every principal ideal of $S$ is a $z$-ideal.
\end{enumerate}
\end{theorem}

\begin{proof}
(1)$\Rightarrow$(2).
Let $I$ be an ideal and let $a\notin I$.
Choose $x\in S$ with $a=a^2x$ and put $e=ax$.
Then $e^2=e$ and $\ideal{a}=\ideal{e}$:
indeed $e=ax\in\ideal{a}$ and $a=ae\in\ideal{e}$.
Since $a\notin I$, also $e\notin I$.
By hypothesis, $e$ has a complement $f$
satisfying $ef=0$ and $e+f=1$.

Consider $J=I+\ideal{f}$.
If $1\in J$, write $1=i+sf$ for some $i\in I$ and $s\in S$,
and multiply by $e$ to get $e=ei\in I$, a contradiction.
Hence $J$ is proper, so some maximal ideal $M$ contains $J$.
Since $f\in M$ and $e+f=1$, we have $e\notin M$, hence $a\notin M$.
Thus every element outside $I$ is excluded by a maximal ideal
containing $I$, proving (2).

(2)$\Rightarrow$(3).
Let $I$ be a strong $z$-ideal, $a\in I$, and $\M(a)=\M(b)$.
Every maximal ideal $M$ containing $I$ contains $a$,
hence $M\in\M(a)=\M(b)$, so $b\in M$.
Intersecting over all such $M$ gives $b\in I$.

(3)$\Rightarrow$(4) is immediate.

(4)$\Rightarrow$(1).
Let $a\in S$.  Since maximal ideals are prime,
Lemma~\ref{loc6}(1) gives $\M(a)=\M(a^2)$.
The ideal $\ideal{a^2}$ is a $z$-ideal by hypothesis, and
$a^2\in\ideal{a^2}$; hence $a\in\ideal{a^2}$,
giving $a=a^2x$ for some $x$.
\end{proof}

\end{document}
